\section[Conclusão]{Conclusão}

Minha experiência como \ac{ages} III no projeto Totem Interativo representou um marco significativo em minha trajetória acadêmica. Pela primeira vez, assumi um papel de liderança técnica, sendo responsável não apenas pelo desenvolvimento, mas também pela orientação dos \ac{ages} I e II. Essa posição trouxe desafios únicos que contribuíram para meu amadurecimento profissional.

O domínio da biblioteca Three.js \cite{threejs} mostrou-se fundamental para o sucesso do projeto. A capacidade de criar mapas 3D interativos, desde versões simples até a complexa representação do Tecnopuc, permitiu que eu assumisse tarefas críticas e garantisse entregas de qualidade, superando as limitações de \textit{hardware} dos totens.

Um dos maiores aprendizados foi sobre a dinâmica de equipe. A parceria com Eduardo Felber Eichner foi crucial para manter o progresso do projeto, mas reconheço que deveria ter sido mais assertivo ao comunicar preocupações sobre a distribuição desigual de responsabilidades entre os \ac{ages} III. Essa experiência destacou a importância da comunicação direta em projetos colaborativos.

Os contratempos técnicos, como problemas com o GitLab e correções de última hora, testaram minha capacidade de resolução rápida de problemas. Manter a calma e encontrar soluções sob pressão foi um aprendizado valioso para minha carreira.

O sucesso da apresentação final, com a validação do \textit{stakeholder} e o interesse de Jorge Audy do Tecnopuc, confirmou a qualidade do trabalho desenvolvido. Ver nossa solução reconhecida como ferramenta com potencial real de aplicação foi extremamente gratificante.

Esta passagem pela \ac{ages} III consolidou conhecimentos técnicos em arquiteturas modernas e habilidades essenciais de liderança e trabalho em equipe. A experiência de revisar código, orientar implementações e tomar decisões arquiteturais preparou-me adequadamente para os desafios do mercado de trabalho, marcando minha evolução de estudante para engenheiro de software.
